% Generated by book/tools/notebook_to_tex.py. Do not edit by hand.

\chapter[Softmax Binary]{sigmoid와 softmax의 동등성}

\markboth{Softmax Binary}{Softmax Binary}

\chaptermeta{04_softmax_binary/04_softmax_binary.ipynb}{https://colab.research.google.com/github/yoon-gu/neuqes-101/blob/master/04_softmax_binary/04_softmax_binary.ipynb}{2차원 softmax 이진 분류와 1차원 sigmoid의 관계}



\index{softmax}

\index{CrossEntropyLoss}

\index{sigmoid}

\index{multinomial}

\index{reparameterization}

\index{소프트맥스}

\index{교차 엔트로피}

\index{원-핫}

\index{리파라미터화}

\index{이진 분류 동등성}



\section{변화추적표}

\begin{adjustbox}{max width=\textwidth}
\begin{tabular}{@{}lllllll@{}}
\toprule
Ch & 모델 & 토크나이저 & 데이터 & Output Head & Activation & Loss \\
\midrule

1 & (TF-IDF) & \inlinecode{TfidfVectorizer()} & Yelp 5,000 & --- & --- & --- \\
2 & \inlinecode{LinearRegression()} & \inlinecode{TfidfVectorizer()} & Yelp (별점 1-5) & (1차원) & 없음 & \inlinecode{MSELoss} \\
3 & \inlinecode{LogisticRegression()} & \inlinecode{TfidfVectorizer()} & Yelp 이진화 & (1차원) & sigmoid & \inlinecode{BCEWithLogitsLoss} \\
\textbf{4 ← 여기} & \inlinecode{LogisticRegression(multi\_class="multinomial")} & \inlinecode{TfidfVectorizer()} & Yelp 이진화 (3장과 동일) & \textbf{(2차원)} & \textbf{softmax} & \textbf{\inlinecode{CrossEntropyLoss}} \\
\bottomrule
\end{tabular}
\end{adjustbox}

전체 19개 장의 표는 \href{https://github.com/yoon-gu/neuqes-101\#장별-변화추적표}{루트 README.md}를 참고하세요.



\section{변경점: 3장 대비}

\begin{adjustbox}{max width=\textwidth}
\begin{tabular}{@{}lll@{}}
\toprule
축 & 3장 & 4장 \\
\midrule

Output Head & (1차원) & \textbf{(2차원)} \\
Activation & sigmoid & \textbf{softmax} \\
Loss & \inlinecode{BCEWithLogitsLoss} & \textbf{\inlinecode{CrossEntropyLoss}} \\
라벨 & int (0/1) & int (0/1) (그대로) \\
데이터 & Yelp 이진화 & Yelp 이진화 (그대로) \\
토크나이저 & TF-IDF & TF-IDF (그대로) \\
\bottomrule
\end{tabular}
\end{adjustbox}

\textbf{왜 같은 데이터에 같은 task인데 따로 장로 다루나?} 두 방식은 출력 차원·activation·loss가 모두 바뀌어 보이지만 \emph{수학적으로 동등} 합니다. 이 동등성을 가장 단순한 sklearn 환경에서 미리 체험해두면, 10장에서 BERT binary가 두 방식 중 어느 쪽을 골라도 같다는 사실을 자연스럽게 받아들일 수 있습니다.

또 K=2를 통과하면 같은 식이 K=5(다음 장)로 자연스럽게 일반화됩니다 --- softmax/CE는 K가 무엇이든 작동합니다.



\section{\texorpdfstring{Loss 함수의 변화 --- \inlinecode{CrossEntropyLoss} 등장}{Loss 함수의 변화 --- CrossEntropyLoss 등장}}

\textbf{Cross Entropy} 는 모델 예측 분포 \(\hat{\mathbf{p}}\) 와 정답 one-hot \(\mathbf{y}\) 사이의 차이를 잽니다.

\begin{equation}
\label{eq:ch04-01}
L = -\frac{1}{N}\sum_{i=1}^{N}\sum_{k=1}^{K} y_{ik} \log \hat p_{ik}
\end{equation}

식~\eqref{eq:ch04-01}은 이 절에서 사용하는 핵심 관계를 정리한 것입니다.

원-핫 정답이라 정답 클래스 항만 살아남습니다.

\begin{equation}
\label{eq:ch04-02}
L = -\frac{1}{N}\sum_{i=1}^{N} \log \hat p_{i,\, y_i}
\end{equation}

식~\eqref{eq:ch04-02}은 이 절에서 사용하는 핵심 관계를 정리한 것입니다.

\textbf{숫자로 감 잡기} (K=2, 정답 \(y = 1\)인 한 샘플):

\begin{adjustbox}{max width=\textwidth}
\begin{tabular}{@{}llll@{}}
\toprule
정답 \(y\) & 예측 분포 \([\hat p_0, \hat p_1]\) & 정답 확률 \(\hat p_1\) & 손실 \(-\log \hat p_1\) \\
\midrule

1 & \texttt{{[}0.1,\ 0.9{]}} & 0.9 & 0.105 \\
1 & \texttt{{[}0.5,\ 0.5{]}} & 0.5 & 0.693 \\
1 & \texttt{{[}0.9,\ 0.1{]}} & 0.1 & \textbf{2.303} \\
\bottomrule
\end{tabular}
\end{adjustbox}

\textbf{3장 BCE 표와 비교해보면} 손실값이 \emph{완전히 같습니다} (0.105 / 0.693 / 2.303). K=2에서 BCE와 CE가 동등하다는 첫 단서.

\begin{lstlisting}
# PyTorch (10장 이후, 방식 B)
criterion = nn.CrossEntropyLoss()
loss = criterion(logits, targets)   # logits: (N, K), targets: (N,) 정수 인덱스

# sklearn (이 장)
LogisticRegression(multi_class="multinomial", max_iter=1000)
\end{lstlisting}



\section{토크나이저 노트}

이 장의 토크나이저는 \textbf{1-3장과 동일한 \inlinecode{TfidfVectorizer}} 입니다. 모델·loss·데이터까지 3장과 같고, 변하는 건 출력 차원과 활성화 함수뿐.

\begin{previewBox}{미리보기}
\textbf{다음 장(5장)}: 같은 TF-IDF, 같은 multinomial LogReg. 변하는 건 데이터가 binary에서 5클래스로, K가 2에서 5로.
\end{previewBox}



\begin{lstlisting}
!pip install -q datasets scikit-learn pandas matplotlib
\end{lstlisting}

\vspace{0.45em}
\noindent\textbf{위 코드 읽기.}\quad \inlinecode{!pip install -q datasets scikit-learn pandas matplotlib}에서는 Colab 실행에 필요한 패키지를 설치합니다.
\par\vspace{0.9em}



\begin{lstlisting}
import warnings
warnings.filterwarnings("ignore", category=FutureWarning)
warnings.filterwarnings("ignore", category=DeprecationWarning)

import numpy as np
import pandas as pd
import matplotlib.pyplot as plt
from datasets import load_dataset
from sklearn.feature_extraction.text import TfidfVectorizer
from sklearn.linear_model import LogisticRegression
from sklearn.model_selection import train_test_split
from sklearn.metrics import accuracy_score, log_loss

plt.rcParams["axes.unicode_minus"] = False

dataset = load_dataset("yelp_review_full")
SAMPLE_SIZE = 5000
ds = dataset["train"].shuffle(seed=42).select(range(SAMPLE_SIZE))
df = ds.to_pandas()
df["star"] = df["label"] + 1   # 0-4 → 1-5
\end{lstlisting}

\vspace{0.45em}
\noindent\textbf{위 코드 읽기.}\quad \inlinecode{numpy, pandas, matplotlib, datasets, scikit-learn} 같은 주요 패키지를 불러옵니다. \inlinecode{warnings.filterwarnings("ignore", category=FutureWarning)}에서는 이후 단계에서 재사용할 중간 값을 계산해 변수에 저장합니다. \inlinecode{warnings.filterwarnings("ignore", category=DeprecationW...}에서는 이후 단계에서 재사용할 중간 값을 계산해 변수에 저장합니다.
\par\vspace{0.9em}



\begin{lstlisting}
# 3장과 동일한 이진화: 별점 3 제외, 4-5 → 1, 1-2 → 0
df_bin = df[df["star"] != 3].copy()
df_bin["y"] = (df_bin["star"] >= 4).astype(int)

X_text_train, X_text_test, y_train, y_test = train_test_split(
    df_bin["text"], df_bin["y"], test_size=0.2, random_state=42, stratify=df_bin["y"],
)

tfidf = TfidfVectorizer(max_features=10000)
X_train = tfidf.fit_transform(X_text_train)
X_test = tfidf.transform(X_text_test)

print(f"X_train: {X_train.shape}, 긍정 비율: {y_train.mean():.1%}")
\end{lstlisting}

\vspace{0.45em}
\noindent\textbf{위 코드 읽기.}\quad \inlinecode{df\_bin = df[df["star"] != 3].copy()}에서는 이후 단계에서 재사용할 중간 값을 계산해 변수에 저장합니다. \inlinecode{df\_bin["y"] = (df\_bin["star"] >= 4).astype(int)}에서는 이후 단계에서 재사용할 중간 값을 계산해 변수에 저장합니다. \inlinecode{X\_text\_train, X\_text\_test, y\_train, y\_test = train\_test...}에서는 학습용 데이터와 평가용 데이터를 분리합니다. \inlinecode{tfidf = TfidfVectorizer(max\_features=10000)}에서는 이후 단계에서 재사용할 중간 값을 계산해 변수에 저장합니다. \inlinecode{X\_train = tfidf.fit\_transform(X\_text\_train)}에서는 텍스트나 라벨을 모델이 처리할 수 있는 수치 표현으로 변환합니다. \inlinecode{X\_test = tfidf.transform(X\_text\_test)}에서는 텍스트나 라벨을 모델이 처리할 수 있는 수치 표현으로 변환합니다.
\par\vspace{0.9em}



\section{실습: 두 방식을 나란히 학습}

\begin{adjustbox}{max width=\textwidth}
\begin{tabular}{@{}lllll@{}}
\toprule
방식 & sklearn 인자 & 출력 차원 & 활성화 & loss \\
\midrule

A (3장 그대로) & \inlinecode{LogisticRegression()} & 1 & sigmoid & BCE \\
B (이 장) & \inlinecode{LogisticRegression(multi\_class="multinomial")} & 2 & softmax & CE \\
\bottomrule
\end{tabular}
\end{adjustbox}



\begin{lstlisting}
model_a = LogisticRegression(max_iter=1000)                           # 방식 A
model_a.fit(X_train, y_train)

model_b = LogisticRegression(multi_class="multinomial", max_iter=1000)  # 방식 B
model_b.fit(X_train, y_train)

acc_a = accuracy_score(y_test, model_a.predict(X_test))
acc_b = accuracy_score(y_test, model_b.predict(X_test))

print(f"방식 A (sigmoid + BCE)  accuracy: {acc_a:.4f}")
print(f"방식 B (softmax + CE)   accuracy: {acc_b:.4f}")
print(f"차이: {abs(acc_a - acc_b):.4f}")
\end{lstlisting}

\vspace{0.45em}
\noindent\textbf{위 코드 읽기.}\quad \inlinecode{model\_a = LogisticRegression(max\_iter=1000)...}에서는 이번 실습에서 관찰할 모델 객체를 정의합니다. \inlinecode{model\_a.fit(X\_train, y\_train)}에서는 학습 데이터로 모델 또는 변환기를 적합합니다. \inlinecode{model\_b = LogisticRegression(multi\_class="multinomial",...}에서는 이번 실습에서 관찰할 모델 객체를 정의합니다. \inlinecode{model\_b.fit(X\_train, y\_train)}에서는 학습 데이터로 모델 또는 변환기를 적합합니다. \inlinecode{acc\_a = accuracy\_score(y\_test, model\_a.predict(X\_test))}에서는 학습된 모델로 최종 예측값을 만듭니다. \inlinecode{acc\_b = accuracy\_score(y\_test, model\_b.predict(X\_test))}에서는 학습된 모델로 최종 예측값을 만듭니다.
\par\vspace{0.9em}



\begin{lstlisting}
proba_a = model_a.predict_proba(X_test)   # (N, 2)
proba_b = model_b.predict_proba(X_test)   # (N, 2)

print(f"방식 A predict_proba shape: {proba_a.shape}")
print(f"방식 B predict_proba shape: {proba_b.shape}")
print(f"  (방식 A도 sklearn이 내부적으로 [P(0), P(1)]을 내줌 — sigmoid 결과를 두 열로 펼친 것)")

p_a, p_b = proba_a[:, 1], proba_b[:, 1]
print(f"\n앞 5개 P(y=1):")
print(f"방식 A: {p_a[:5].round(4)}")
print(f"방식 B: {p_b[:5].round(4)}")
print(f"\n전체 max 차이: {np.abs(p_a - p_b).max():.4f}")
print(f"전체 mean 차이: {np.abs(p_a - p_b).mean():.4f}")
\end{lstlisting}

\vspace{0.45em}
\noindent\textbf{위 코드 읽기.}\quad \inlinecode{proba\_a = model\_a.predict\_proba(X\_test)   \# (N, 2)}에서는 클래스별 예측 확률을 계산합니다. \inlinecode{proba\_b = model\_b.predict\_proba(X\_test)   \# (N, 2)}에서는 클래스별 예측 확률을 계산합니다. \inlinecode{p\_a, p\_b = proba\_a[:, 1], proba\_b[:, 1]}에서는 이후 단계에서 재사용할 중간 값을 계산해 변수에 저장합니다.
\par\vspace{0.9em}



\section{해부: 수학적 동등성}

방식 B는 \emph{개념적으로} logit을 두 개 (\(z_0, z_1\)) 학습합니다. softmax의 두 번째 성분을 풀어보면:

\begin{equation}
\label{eq:ch04-03}
\text{softmax}([z_0, z_1])_1 = \frac{e^{z_1}}{e^{z_0} + e^{z_1}} = \frac{1}{1 + e^{-(z_1 - z_0)}} = \sigma(z_1 - z_0)
\end{equation}

식~\eqref{eq:ch04-03}은 2차원 softmax가 logit 차이에 대한 sigmoid로 표현됨을 보여줍니다.

즉 두 logit에서 \textbf{의미 있는 정보는 \(z_1 - z_0\) 뿐} 입니다 --- 두 logit에 같은 상수를 더해도 softmax 결과가 안 바뀌니까요 (\(e^{z+c}/\sum e^{z+c} = e^{z}/\sum e^{z}\)). softmax+2는 sigmoid+1의 \textbf{리파라미터화} 일 뿐입니다.

CE 쪽도 K=2에서 BCE와 같은 식이 됩니다 (one-hot이라 \(y_1 = y\), \(y_0 = 1-y\) 대입):

\begin{equation}
\label{eq:ch04-04}
\text{CE} = -[y_1 \log \hat p_1 + y_0 \log \hat p_0] = -[y \log \hat p_1 + (1-y)\log(1 - \hat p_1)] = \text{BCE}
\end{equation}

식~\eqref{eq:ch04-04}은 정답 라벨과 예측 확률 사이의 Binary Cross-Entropy를 정의합니다.

확률·loss가 같으니 학습된 결정 경계도, gradient도 같습니다.

먼저 식이 정말 일치하는지 임의의 logit 쌍으로 직접 확인합니다.



\begin{lstlisting}
# 임의의 logit 쌍 4개를 만들어 softmax([z0,z1])_1 == sigmoid(z1 - z0) 인지 확인
z0_arr = np.array([-2.0, 0.0, 1.5, 3.0])
z1_arr = np.array([ 1.0, 0.5, -0.5, 2.0])

softmax_p1  = np.exp(z1_arr) / (np.exp(z0_arr) + np.exp(z1_arr))
sigmoid_diff = 1.0 / (1.0 + np.exp(-(z1_arr - z0_arr)))

print(f"{'z_0':>6} {'z_1':>6}    {'softmax([z0,z1])_1':>22}    {'sigmoid(z1-z0)':>16}")
print("-" * 60)
for i in range(len(z0_arr)):
    print(f"{z0_arr[i]:>6.1f} {z1_arr[i]:>6.1f}    {softmax_p1[i]:>22.8f}    {sigmoid_diff[i]:>16.8f}")

print(f"\nmax 차이: {np.abs(softmax_p1 - sigmoid_diff).max():.2e}  (수치적 오차 수준)")
\end{lstlisting}

\vspace{0.45em}
\noindent\textbf{위 코드 읽기.}\quad \inlinecode{z0\_arr = np.array([-2.0, 0.0, 1.5, 3.0])}에서는 이후 단계에서 재사용할 중간 값을 계산해 변수에 저장합니다. \inlinecode{z1\_arr = np.array([ 1.0, 0.5, -0.5, 2.0])}에서는 이후 단계에서 재사용할 중간 값을 계산해 변수에 저장합니다. \inlinecode{softmax\_p1  = np.exp(z1\_arr) / (np.exp(z0\_arr) + np.exp...}에서는 이후 단계에서 재사용할 중간 값을 계산해 변수에 저장합니다. \inlinecode{sigmoid\_diff = 1.0 / (1.0 + np.exp(-(z1\_arr - z0\_arr)))}에서는 이후 단계에서 재사용할 중간 값을 계산해 변수에 저장합니다. \inlinecode{for i in range(len(z0\_arr)):}에서는 앞 단계에서 만든 값을 바탕으로 다음 계산을 수행합니다.
\par\vspace{0.9em}



\section{\texorpdfstring{변형: sklearn은 왜 K=2 multinomial에서 \inlinecode{(2,\ V)} coef를 안 만드나?}{변형: sklearn은 왜 K=2 multinomial에서 (2, V) coef를 안 만드나?}}

위 동등성 덕분에 K=2에서 두 logit 중 하나는 잉여입니다. sklearn은 이 사실을 알고 \textbf{K=2 multinomial을 자동으로 binary form으로 collapse} 시킵니다 --- \inlinecode{coef\_} 를 \inlinecode{(2,\ V)} 가 아니라 \inlinecode{(1,\ V)} 로만 저장합니다. 두 방식이 그래서 사실상 같은 모델이 되어 \inlinecode{predict\_proba}도 거의 일치하는 거였죠.

직접 두 모델의 \inlinecode{coef\_} 모양을 확인합니다.



\begin{lstlisting}
print(f"방식 A coef_ shape:      {model_a.coef_.shape}")
print(f"방식 B coef_ shape:      {model_b.coef_.shape}")
print(f"방식 A intercept_ shape: {model_a.intercept_.shape}")
print(f"방식 B intercept_ shape: {model_b.intercept_.shape}")
print()
print("→ 둘 다 (1, V) — sklearn이 K=2 multinomial을 binary form으로 collapse")
print()
print(f"두 coef_ 의 max 차이:      {np.abs(model_a.coef_ - model_b.coef_).max():.2e}")
print(f"두 intercept_ 의 max 차이: {np.abs(model_a.intercept_ - model_b.intercept_).max():.2e}")
print()
print("(미세한 차이는 solver 수렴 기준의 차이일 뿐, 본질적으로 같은 모델)")
print()
print("진짜 (2, V) 두 logit head는 PyTorch가 '직접' 만들어주는 10장 BERT binary에서 등장합니다.")
\end{lstlisting}



\section{이 장에 등장한 라이브러리}

\begin{adjustbox}{max width=\textwidth}
\begin{tabular}{@{}lll@{}}
\toprule
이름 & 한 줄 설명 & 다음 장에서 \\
\midrule

\inlinecode{LogisticRegression(multi\_class="multinomial")} & softmax + CE 다중 분류 (이 장엔 K=2) & 5장에서 K=5로 확장, 11장 BERT multi-class에서 같은 패러다임 \\
\inlinecode{sklearn.metrics.log\_loss} & CE/BCE 평가 함수 (multi-class 호환) & --- \\
\bottomrule
\end{tabular}
\end{adjustbox}



\section{체크포인트 질문}

\begin{enumerate}
\def\labelenumi{\arabic{enumi}.}
\tightlist
\item
  softmax 함수 정의를 적고, \(\text{softmax}([z_0, z_1])_1\) 이 \(\sigma(z_1 - z_0)\) 와 같음을 증명해보세요.
\item
  Cross Entropy를 K=2에 적용하면 정확히 BCE가 되는 과정을 식으로 보일 수 있나요? (\(y_1 = y\), \(y_0 = 1-y\) 대입)
\item
  방식 B의 두 coefficient 벡터 사이에 어떤 관계가 학습되는 경향이 있나요? 그 이유는?
\item
  같은 binary 데이터에 두 방식의 accuracy가 거의 같다면, 실무에서 어느 쪽을 택해야 하나요?
\end{enumerate}



\section{FAQ}

\begin{faqBox}{Q1. (이론) K=2일 때 sigmoid+BCE와 softmax+CE가 정확히 동등하다는 걸 식으로 어떻게 보이나요?}

두 가지를 따로 보여야 합니다.

\textbf{확률 동등성:}

\begin{equation}
\label{eq:ch04-05}
\text{softmax}([z_0, z_1])_1 = \frac{e^{z_1}}{e^{z_0}+e^{z_1}} = \frac{1}{1+e^{-(z_1-z_0)}} = \sigma(z_1 - z_0)
\end{equation}

식~\eqref{eq:ch04-05}은 2차원 softmax가 logit 차이에 대한 sigmoid로 표현됨을 보여줍니다.

방식 A의 logit을 \(z = z_1 - z_0\) 로 두면 두 모델의 P(y=1)이 일치합니다.

\textbf{Loss 동등성} (\(K=2\), one-hot 정답에서 \(y_1 = y\), \(y_0 = 1-y\)):

\begin{equation}
\label{eq:ch04-06}
\text{CE} = -\sum_{k=0}^{1} y_k \log \hat p_k = -[y \log \hat p_1 + (1-y) \log \hat p_0] = -[y \log \hat p_1 + (1-y) \log(1 - \hat p_1)] = \text{BCE}
\end{equation}

식~\eqref{eq:ch04-06}은 정답 라벨과 예측 확률 사이의 Binary Cross-Entropy를 정의합니다.

확률과 loss가 모두 같으니 학습된 결정 경계도 같고, gradient도 같습니다.

\end{faqBox}

\begin{faqBox}{Q2. (실무) 실제로 둘 중 어느 방식이 더 널리 쓰이나요?}

두 가지 관행이 공존합니다.

\begin{itemize}
\tightlist
\item
  \textbf{sigmoid+BCE (방식 A, num\_labels=1)}: sklearn 기본, 통계학·의학 분야 표준. 출력 1개라 "확률 하나"라는 해석이 단순.
\item
  \textbf{softmax+CE (방식 B, num\_labels=2)}: BERT/PyTorch 기본, 딥러닝 표준. 다중 클래스로 일반화하기 자연스럽고 라이브러리 코드가 단순(같은 헤드/같은 loss로 K가 2이든 N이든 호환).
\end{itemize}

이 커리큘럼의 BERT 장(9-13장)는 방식 B가 기본이라 이 장에서 미리 익숙해지는 게 의미 있습니다. 10장에서 두 방식을 BERT로 다시 비교합니다.

\end{faqBox}

\begin{faqBox}{Q3. (이론) softmax 합=1 제약은 어디서 오나요?}

수식 자체가 정규화를 강제합니다.

\begin{equation}
\label{eq:ch04-07}
\text{softmax}(z)_k = \frac{e^{z_k}}{\sum_{j=1}^{K} e^{z_j}}
\end{equation}

식~\eqref{eq:ch04-07}은 logit 벡터를 확률 분포로 바꾸는 softmax의 정의입니다.

분모가 모든 클래스 \(e^{z_j}\)의 합이라 분자들이 그 합을 정확히 분할 --- 합이 1.

\textbf{왜 이런 구조?} 모델 출력을 "확률 분포"로 해석하고 싶어서입니다. 확률 분포는 정의상 \(\sum_k p_k = 1\) 이고 각 항이 {[}0, 1{]}. softmax는 임의 실수 logit 벡터를 그런 분포로 보내는 가장 자연스러운 변환 중 하나(지수 함수 = 단조증가 + 양수 보장 → 정규화).

\end{faqBox}

\begin{faqBox}{Q4. (실무) sklearn에서 binary에 \inlinecode{multi\_class="multinomial"} 을 줘도 \inlinecode{coef\_.shape}가 \inlinecode{(1,\ V)} 인 이유는?}

위 본문에서 확인했듯 K=2 softmax는 두 logit 중 하나가 잉여(redundant)입니다 --- \(z_1 - z_0\) 만 의미가 있어요. sklearn은 이걸 알고 \textbf{K=2 multinomial을 자동으로 binary form으로 collapse} 시킵니다. 그래서 \inlinecode{coef\_} 가 \inlinecode{(2,\ V)} 가 아니라 \inlinecode{(1,\ V)}, \inlinecode{intercept\_} 도 \inlinecode{(1,)} 로 저장됩니다.

\begin{lstlisting}
LogisticRegression(multi_class="multinomial").fit(X, y_binary).coef_.shape  # (1, V)
LogisticRegression(multi_class="multinomial").fit(X, y_3class).coef_.shape  # (3, V) — K≥3에선 (K, V)
\end{lstlisting}

그래서 방식 A와 방식 B가 sklearn 안에서는 사실상 같은 모델이고, predict\_proba도 미세한 수치 오차 빼고 일치합니다. 진짜 \emph{두 별개의 logit head} 가 살아 있는 형태는 프레임워크가 collapse하지 않는 환경 --- PyTorch에서 \inlinecode{nn.Linear(H,\ 2)} 를 직접 만들 때 --- 비로소 등장합니다 (10장).

\end{faqBox}

\begin{faqBox}{Q5. (이론) sklearn \inlinecode{multi\_class} 인자의 의미와 자동 선택 기준은?}

\begin{itemize}
\tightlist
\item
  \inlinecode{"multinomial"}: 모든 클래스를 한 번에 학습 (softmax). 이 장.
\item
  \inlinecode{"ovr"} (One-vs-Rest): K개 독립 binary 분류기. 6장 multi-label로 가는 다리.
\item
  \inlinecode{"auto"} (기본값): solver와 데이터에 맞춰 자동 --- 클래스 2개면 binary, 3개 이상이면 multinomial(\inlinecode{lbfgs}, \inlinecode{newton-cg}, \inlinecode{sag}, \inlinecode{saga} solver) 또는 OvR(\inlinecode{liblinear}).
\end{itemize}

sklearn 1.5부터 \inlinecode{multi\_class} 자체가 deprecated이고 multinomial이 기본 동작이 됩니다. 호환성을 위해 명시할 때만 \inlinecode{multi\_class="multinomial"} 처럼 적습니다.

\end{faqBox}

\begin{faqBox}{Q6. (실무) Hugging Face \inlinecode{Trainer}도 두 방식이 가능한가요?}

가능하고, \inlinecode{AutoModelForSequenceClassification} 인자 한 줄 차이입니다.

\begin{lstlisting}
# 방식 A: num_labels=1, BCE 자동 적용
AutoModelForSequenceClassification.from_pretrained(
    "distilbert-base-uncased",
    num_labels=1,
    problem_type="single_label_classification",  # 또는 자동
)

# 방식 B: num_labels=2, CE 자동 적용 (BERT 표준)
AutoModelForSequenceClassification.from_pretrained(
    "distilbert-base-uncased",
    num_labels=2,
)
\end{lstlisting}

이 커리큘럼의 10장이 두 방식을 BERT로 다시 비교하는 장입니다. 이 장에서 익힌 동등성이 그 비교의 출발점이 됩니다.
\end{faqBox}



\section{오류 실험 (선택)}

다음 코드는 sklearn 없이 numpy만으로 softmax의 \textbf{shift invariance} 를 보여줍니다 --- 모든 logit에 같은 상수를 더해도 출력 분포가 변하지 않는다는 성질입니다. 이게 K=2일 때 sklearn이 collapse를 하는 \emph{정확한 이유} 입니다.

\begin{lstlisting}
z = np.array([1.0, 2.0])

p_original  = np.exp(z)         / np.sum(np.exp(z))
p_shift5    = np.exp(z + 5)     / np.sum(np.exp(z + 5))
p_zero_z0   = np.exp(z - z[0])  / np.sum(np.exp(z - z[0]))   # z_0 = 0 으로 정규화한 형태

print(f"softmax([1, 2]):    {p_original.round(6)}")
print(f"softmax([6, 7]):    {p_shift5.round(6)}")
print(f"softmax([0, 1]):    {p_zero_z0.round(6)}")
\end{lstlisting}

힌트: 세 결과가 모두 같습니다. 두 logit 중 \emph{하나는 자유롭게 정할 수 있고} 의미 있는 정보는 차이 \(z_1 - z_0\) 뿐 --- 만약 \(z_0 = 0\) 으로 고정하면 \(z_1\) 만 학습하면 됩니다. 이게 sklearn이 K=2 multinomial을 binary form으로 collapse하는 형식적 근거이고, K=2 softmax가 sigmoid의 리파라미터화에 불과한 이유입니다.



\begin{previewBox}{미리보기: 다음 장}

\textbf{5장. sklearn Multi-class --- K=5로 진짜 일반화}

\begin{itemize}
\tightlist
\item
  같은 multinomial LogReg를 별점 1-5(5클래스)로 그대로 확장
\item
  수식·코드 변화 거의 없음 --- softmax/CE는 K가 무엇이든 같은 형태
\item
  5×5 confusion matrix가 대각선 근처에 몰리는 ordinal 흔적 관찰
\item
  multinomial vs OvR 비교 (6장 multi-label로 가는 다리)
\end{itemize}
\end{previewBox}
